\section{Simultaneous cost-and-cut rounding}
\label{sec:fractional-rounding}

The all-cut theorem also permits control of one nonnegative cost
functional. The reduction combines the truncation and rounding in
\eqref{eq:ta-binary-reduction} with repeated deletion and selection of
the least-cost tree, as in Oveis Gharan and Saberi
\cite[full preprint, Proposition~5.1]{OveisGharanSaberi2011}.
The rational compression below keeps the auxiliary graph polynomial
in the input bit length.

\begin{corollary}[Thin and cheap rounding of fractional cuts]
\label{cor:ta-fractional-rounding}
Fix once and for all an integer $B\ge\max\{1,C\}$, where $C$ is the
constant in Theorem~\ref{thm:tt-algorithm}. Let $G=(V,E)$ be a finite undirected
graph on $n\ge2$ vertices, given by an explicit list of edge records;
in this corollary, loops and parallel records are allowed. Given
nonnegative rational numbers $x_e,c_e$ in binary, suppose that
\[
 x(\delta_G(S)):=\sum_{e\in\delta_G(S)}x_e\ge1
       \qquad(\varnothing\ne S\subsetneq V).
\]
There is a deterministic algorithm, polynomial in the total binary
input length, that returns a spanning tree
$T\subseteq\{e\in E:x_e>0\}$ satisfying simultaneously
\begin{align}
 |\delta_T(S)|&\le24B\,x(\delta_G(S))
       &&(\varnothing\ne S\subsetneq V),
       \label{eq:ta-fractional-cuts}\\
 c(T):=\sum_{e\in T}c_e&\le24B\sum_{e\in E}c_ex_e.
       \label{eq:ta-fractional-cost}
\end{align}
The cut inequalities are a promise on the input; no metric assumption
on the costs is required.
\end{corollary}

\begin{proof}
Discard loops, which affect no cut and can only reduce the nonnegative
cost sum. For each unordered nonloop pair $p$, let $E_p$ be its set of
edge records and put $X_p=\sum_{e\in E_p}x_e$. Discard pairs with $X_p=0$;
otherwise choose an original edge $e_p$ with $x_{e_p}>0$ of minimum
cost among such edges, and put $d_p=c_{e_p}$. This preserves every
cut mass and gives
\[
 \sum_p d_pX_p\le\sum_{e\in E}c_ex_e.
\]
Set $K=n^2$, $y_p=\min\{X_p,1\}$, and
$a_p=\lfloor2K y_p\rfloor$. Form an explicit loopless multigraph $A$
with $a_p$ labelled copies of pair $p$, each of cost $d_p$.
Every cut has $y$-mass at least one: if a crossing coordinate was
truncated, that coordinate alone contributes one, and otherwise the
cut mass is unchanged. A cut crosses at most $n^2/4$ unordered pairs,
so rounding down loses at most $n^2/4$ from its scaled mass. Hence
\begin{align}
 |\delta_A(S)|&\ge2K-n^2/4\ge K,
       \label{eq:ta-fractional-lower}\\
 |\delta_A(S)|&\le2K\,x(\delta_G(S)),
       \qquad
 c(A):=\sum_p a_pd_p\le2K\sum_{e\in E}c_ex_e.
       \label{eq:ta-fractional-upper}
\end{align}
Each $a_p\le2n^2$, and fewer than $n^2/2$ pairs occur, so
$|E(A)|<n^4$. In particular, $A$ is connected.

If $K<12B$, select any spanning tree of $A$. Its load on each cut is
at most that of $A$, and its cost is at most $c(A)$ by nonnegativity.
Equation~\eqref{eq:ta-fractional-upper} and $2K<24B$ therefore give
both asserted bounds, including when the cost sum is zero.

Suppose now that $K\ge12B$, and set
$k=\lfloor K/2\rfloor$ and $m=\lfloor K/(6B)\rfloor$.
Starting with $R_0=A$, apply Theorem~\ref{thm:tt-algorithm} with
connectivity parameter $k$ to $R_i$, obtain a tree $T_i$, and delete
its labelled edges to form $R_{i+1}$, for $0\le i<m$.
To verify that every call is legal, inductively we have, for all cuts
and $0\le i\le m$,
\begin{equation}
 |\delta_{R_i}(S)|\ge
       (1-3Bi/K)|\delta_A(S)|\ge\tfrac12|\delta_A(S)|.
 \label{eq:ta-fractional-residual}
\end{equation}
Indeed, the second inequality follows from $m\le K/(6B)$.
Together with \eqref{eq:ta-fractional-lower}, it makes $R_i$
$k$-edge-connected before its call. Since $k\ge K/3$, the returned tree
obeys
\[
 |\delta_{T_i}(S)|\le\frac Ck|\delta_{R_i}(S)|
       \le\frac{3B}{K}|\delta_A(S)|,
\]
which proves the next instance of
\eqref{eq:ta-fractional-residual}. Equation~\eqref{eq:ta-fractional-upper}
then makes every $T_i$ at most $6B$-thin relative to $x$.
The trees are edge-disjoint in $A$, and
$m\ge K/(12B)$ because $K/(6B)\ge2$. Nonnegative costs therefore give
\[
 \min_{0\le i<m}c(T_i)
       \le\frac{c(A)}m
       \le24B\sum_{e\in E}c_ex_e.
\]
Select a least-cost $T_i$.

In either case, a selected tree uses at most one copy of any parallel
pair. Replace that copy by $e_p$. The resulting original edges still
form a spanning tree, have the same cut loads and cost, and all have
positive $x_e$. Different extracted trees may lift to the same original
edge, which is harmless because only one tree is output.

All pair sums, truncations, floors, and cost comparisons use exact
rational or integer arithmetic of polynomial bit length; no common
denominator is expanded into edge copies. The positive support is
connected by the cut promise, so it contains at least $n-1$ distinct
nonloop pairs. Thus even if $n$ is encoded in binary, it is bounded by
the length of the explicit input. There are $O(n^2)$ calls on graphs
with fewer than $n^4$ explicit edges and an $O(\log n)$-bit parameter
$k$. Deleting labelled copies, summing and comparing tree costs, and
lifting the selected tree are polynomial-bit operations. Very small
positive coordinates may round to zero, but
\eqref{eq:ta-fractional-lower} preserves connectivity. No division by
$\sum_e c_ex_e$ occurs; if this sum is zero, $c(A)=0$ and the output
cost is zero. This proves the stated bit bound and covers zero values.
\end{proof}

\paragraph{The historical thin-and-cheap formulation.}
The author full manuscript of Asadpour, Goemans, M\k{a}dry, Oveis
Gharan, and Saberi defines an $(\alpha,s)$-thin tree for a point $z$
in a spanning-tree polytope by the all-cut bound
$|\delta_T(S)|\le\alpha z(\delta(S))$ and the cost bound
$c(T)\le s c(z)$, with $\alpha\ge1$. Its Conjecture~6.1 asks for
absolute constants and a polynomial-time construction for every
feasible Held--Karp solution
\cite[author full manuscript, Definition~4.1 and Conjecture~6.1]{AsadpourEtAl2017}.

For rational data on $n\ge2$ vertices in that paper's directed-metric
setting, let $h$ be a feasible Held--Karp vector and use its normalized
symmetrization
\[
 z_{\{u,v\}}=\frac{n-1}{n}(h_{uv}+h_{vu})
\]
from equation~(3.5) of the full manuscript. The feasibility calculation
in the proof of its Lemma~3.1 puts $z$ in the corresponding spanning-tree
polytope \cite[author full manuscript, equation~(3.5) and proof of Lemma~3.1]{AsadpourEtAl2017}.
In particular, flow balance and the cut constraint give
\[
 z(\delta(S))=\frac{2(n-1)}n h(\delta^+(S))
       \ge\frac{2(n-1)}n\ge1.
\]
For each support pair $p=\{u,v\}$, set
$d_p=\min\{c_{uv},c_{vu}\}$, the cost of the cheaper of both
orientations. Its induced fractional cost satisfies
\[
 c(z)=\sum_p d_pz_p\le\frac{n-1}{n}\sum_a c_a h_a.
\]
Apply Corollary~\ref{cor:ta-fractional-rounding} with $x=z$ and costs
$d_p$. With the same cost function on the tree and the fractional
point, the corollary gives $c(T)\le24B c(z)$ and the required all-cut
bound. It therefore supplies the two $(24B,24B)$ inequalities on the
undirected support, with deterministic polynomial-bit construction
for rationally encoded data. The cheaper orientation of a selected
pair need not have positive $h$-value; the support requirement is
$z_p>0$. More generally,
every rational point in a spanning-tree polytope satisfies the cut
promise, since each constituent spanning tree meets every nontrivial
cut. The bit guarantee here is not a claim for arbitrary real-oracle
inputs.

The earlier SODA version instead fixes an optimum Held--Karp solution
and defines its cost bound as $c(T)\le s\,OPT_{\rm HK}$ in
Definition~5.1; its Theorem~6.1 is the tour-augmentation theorem, not
the full-manuscript conjecture
\cite[Definition~5.1 and Theorem~6.1]{AsadpourEtAl2010}.
Constant-factor approximation for ATSP was already established
independently by Svensson, Tarnawski, and V\'egh and improved by Traub
and Vygen \cite{SvenssonTarnawskiVegh2020,TraubVygen2022}.
The consequence proved here concerns the underlying simultaneous
rounding formulation; no new approximation-order claim is made.
